\CWHeader{Laboratory work No.~1}

\begin{table}[!ht]
\centering
\begin{tabular}{|l|l|}
\hline
\rowcolor{lightgray}
\multicolumn{2}{|c|}{J. 4-1} \\
\hline
Time limit & 3 seconds \\
\hline
Memory limit & 300 MB \\
\hline
Input & standard input or \texttt{input.txt} \\
\hline
Output & standard output or \texttt{output.txt} \\
\hline
\end{tabular}
\end{table}

{\sloppy
\CWProblem{
Implement a program that reads key-value pairs, sorts them by key in ascending order
using a linear-time sorting algorithm, and outputs the sorted sequence. The variant is
defined by the key type (and the corresponding sorting method) and the value type.

{\bfseries Sorting method:} radix sort.

{\bfseries Key type:} dates in the format \texttt{DD.MM.YYYY}, e.g. \texttt{1.1.1},
\texttt{1.9.2009}, \texttt{01.09.2009}, \texttt{31.12.2009}.

{\bfseries Value type:} fixed-length strings of 64 characters. Input strings may be
shorter; in that case they are padded with null characters up to 64, which are not printed.
}
\par}

\subsection*{Input format}

Each non-empty line of the input contains a key-value pair: the key as specified above,
followed by a tab character, followed by the corresponding value.

\subsection*{Output format}

The output contains the same lines as the input (excluding empty ones), in sorted order.

\subsection*{Example}

Input:

{\footnotesize
\begin{alltt}
1.1.1       n399tann9nnt3ttnaaan9nann93na9t3a3t9999na3aan9antt3tn93aat3naatt
01.02.2008  n399tann9nnt3ttnaaan9nann93na9t3a3t9999na3aan9antt3tn93aat3naat
1.1.1       n399tann9nnt3ttnaaan9nann93na9t3a3t9999na3aan9antt3tn93aat3naa
01.02.2008  n399tann9nnt3ttnaaan9nann93na9t3a3t9999na3aan9antt3tn93aat3na
\end{alltt}
}

Output:

{\footnotesize
\begin{alltt}
1.1.1       n399tann9nnt3ttnaaan9nann93na9t3a3t9999na3aan9antt3tn93aat3naatt
1.1.1       n399tann9nnt3ttnaaan9nann93na9t3a3t9999na3aan9antt3tn93aat3naa
01.02.2008  n399tann9nnt3ttnaaan9nann93na9t3a3t9999na3aan9antt3tn93aat3naat
01.02.2008  n399tann9nnt3ttnaaan9nann93na9t3a3t9999na3aan9antt3tn93aat3na
\end{alltt}
}

\pagebreak
